\documentclass[10pt,a4paper]{amsart}
\usepackage{amsmath,amssymb,amsthm,graphicx}
\usepackage{verbatim}
\newenvironment{reference}{\comment}{\endcomment}
%\newenvironment{reference}{}{}
\newenvironment{slogan}{\comment}{\endcomment}
\newenvironment{history}{\comment}{\endcomment}
% For commutative diagrams we use Xy-pic
\usepackage[all]{xy}

% We use 2cell for 2-commutative diagrams.
\xyoption{2cell}
% Theorem environments.
%
\theoremstyle{plain}
\newtheorem{theorem}[subsection]{Theorem}
\newtheorem{proposition}[subsection]{Proposition}
\newtheorem{lemma}[subsection]{Lemma}

\theoremstyle{definition}
\newtheorem{definition}[subsection]{Definition}
\newtheorem{example}[subsection]{Example}
\newtheorem{exercise}[subsection]{Exercise}
\newtheorem{situation}[subsection]{Situation}

\theoremstyle{remark}
\newtheorem{remark}[subsection]{Remark}
\newtheorem{remarks}[subsection]{Remarks}

\numberwithin{equation}{subsection}

% Macros
%
\def\lim{\mathop{\mathrm{lim}}\nolimits}
\def\colim{\mathop{\mathrm{colim}}\nolimits}
\def\Spec{\mathop{\mathrm{Spec}}}
\def\Hom{\mathop{\mathrm{Hom}}\nolimits}
\def\Ext{\mathop{\mathrm{Ext}}\nolimits}
\def\SheafHom{\mathop{\mathcal{H}\!\mathit{om}}\nolimits}
\def\SheafExt{\mathop{\mathcal{E}\!\mathit{xt}}\nolimits}
\def\Sch{\mathit{Sch}}
\def\Mor{\mathop{\mathrm{Mor}}\nolimits}
\def\Ob{\mathop{\mathrm{Ob}}\nolimits}
\def\Sh{\mathop{\mathit{Sh}}\nolimits}
\def\NL{\mathop{N\!L}\nolimits}
\def\CH{\mathop{\mathrm{CH}}\nolimits}
\def\proetale{{pro\text{-}\acute{e}tale}}
\def\etale{{\acute{e}tale}}
\def\QCoh{\mathit{QCoh}}
\def\Ker{\mathop{\mathrm{Ker}}}
\def\Im{\mathop{\mathrm{Im}}}
\def\Coker{\mathop{\mathrm{Coker}}}
\def\Coim{\mathop{\mathrm{Coim}}}

% Boxtimes
%
\DeclareMathSymbol{\boxtimes}{\mathbin}{AMSa}{"02}

%
% Macros for moduli stacks/spaces
%
\def\QCohstack{\mathcal{QC}\!\mathit{oh}}
\def\Cohstack{\mathcal{C}\!\mathit{oh}}
\def\Spacesstack{\mathcal{S}\!\mathit{paces}}
\def\Quotfunctor{\mathrm{Quot}}
\def\Hilbfunctor{\mathrm{Hilb}}
\def\Curvesstack{\mathcal{C}\!\mathit{urves}}
\def\Polarizedstack{\mathcal{P}\!\mathit{olarized}}
\def\Complexesstack{\mathcal{C}\!\mathit{omplexes}}
% \Pic is the operator that assigns to X its picard group, usage \Pic(X)
% \Picardstack_{X/B} denotes the Picard stack of X over B
% \Picardfunctor_{X/B} denotes the Picard functor of X over B
\def\Pic{\mathop{\mathrm{Pic}}\nolimits}
\def\Picardstack{\mathcal{P}\!\mathit{ic}}
\def\Picardfunctor{\mathrm{Pic}}
\def\Deformationcategory{\mathcal{D}\!\mathit{ef}}
\usepackage{cleveref}
\crefname{originallemma}{original lemma}{original lemmas}
\Crefname{originallemma}{Original lemma}{Original lemmas}
\makeatletter
\newlabel{divisors-lemma-pullback-effective-Cartier-defined}{{02OO}{source}{}{sourcepointer-1}{}}
\newlabel{divisors-lemma-pullback-effective-Cartier-defined@cref}{{[originallemma][0][]02OO}{[1][0][]source}{}{}{}}
\makeatother
\crefname{originallemma}{original lemma}{original lemmas}
\Crefname{originallemma}{Original lemma}{Original lemmas}
\makeatletter
\newlabel{lemma-proper-pushforward-rational-equivalence}{{02S2}{source}{}{sourcepointer-2}{}}
\newlabel{lemma-proper-pushforward-rational-equivalence@cref}{{[originallemma][0][]02S2}{[1][0][]source}{}{}{}}
\makeatother
\crefname{originallemma}{original lemma}{original lemmas}
\Crefname{originallemma}{Original lemma}{Original lemmas}
\makeatletter
\newlabel{lemma-rational-function}{{02RQ}{source}{}{sourcepointer-3}{}}
\newlabel{lemma-rational-function@cref}{{[originallemma][0][]02RQ}{[1][0][]source}{}{}{}}
\makeatother
\crefname{originalreference}{original reference}{original references}
\Crefname{originalreference}{Original reference}{Original references}
\makeatletter
\newlabel{situation-setup}{{02QL}{source}{}{sourcepointer-4}{}}
\newlabel{situation-setup@cref}{{[originalreference][0][]02QL}{[1][0][]source}{}{}{}}
\makeatother
\title{Rational equivalence on a product with the projective line: the two fibres of an integral closed subscheme and the generation of rationally trivial cycles}
\author{The Stacks Project Authors}
\date{stacks-project:chow.tex:local-snapshot-20261012; GNU FDL 1.2 or later}
\begin{document}
\maketitle
\noindent\textit{Editorial note.} All delivered mathematical text is the author's own, copied from the pinned Stacks Project chapter file chow.tex; no statement or proof step was rewritten, reordered or removed.\par\medskip
\noindent\textit{Editorial note.} Preamble provenance: the pinned chapter file opens with \textbackslash{}input\{preamble\} and carries no preamble text of its own, so the project's single Stacks-wide preamble staged at sources/stacks-chow/preamble.tex is excerpted verbatim as three spans, lines 8-12 (the verbatim package together with the project's own reference, slogan and history display-environment definitions, including the author's commented-out older form of the reference environment), lines 14-18 (Xy-pic setup) and lines 154-217 (the native theorem declarations and mathematical macros). The lines 8-12 span is required because a counted statement or its surrounding text may use those display environments, which the project typesets as comments. Nothing was added, removed or redefined; the wrapper class and the amsmath/amssymb/amsthm/graphicx base packages are supplied by the assembler.\par\medskip
\noindent\textit{Editorial note.} Same-chapter pointers are located at their \textbackslash{}label lines inside the pinned chapter file; cross-chapter pointers are located in the project's own tag list staged in the snapshot (sources/stacks-chow/tags.txt), at the verbatim lines carrying their canonical tags.\par\medskip
\noindent\textit{Allowed standard prerequisites (editorial).}\begin{itemize}
\item The two fibres of a closed subscheme of the product as effective Cartier divisors. Divisors, label divisors-lemma-pullback-effective-Cartier-defined, invoked at proof line 3364. Gives part (1): because the generic point of W is not mapped into either of the two standard sections, the two fibres are effective Cartier divisors of W.
\item The graph construction attached to a rational function on an integral scheme. Same chapter chow.tex, label lemma-rational-function, invoked at proof lines 3375 and 3397. Supplies both the rational equivalence of the two fibre classes on W for part (2) and the closed subscheme W\_i inside V\_i times the projective line whose two fibre classes realise the divisor of f\_i in part (4).
\item A proper morphism preserves rational equivalence. Same chapter chow.tex, label lemma-proper-pushforward-rational-equivalence, invoked at proof line 3377; the definition of rational equivalence is definition-rational-equivalence (used at proof lines 3355 and 3386). Moves the rational equivalence from W to X, since the projection is proper, which is what turns the equivalence on the graph into the statement about the fibres in X.
\end{itemize}\medskip
\ifdefined\hypertarget\hypertarget{sourcepointer-1}{}\fi
\noindent\textit{Source reference.} Original lemma 02OO: sources/stacks-chow/tags.txt, source lines 3523--3523. Navigation pointer for \textbackslash{}ref\{divisors-lemma-pullback-effective-Cartier-defined\} used inside the counted span (line 3364); the argument that consumes it is retained in the span.\par\medskip
\ifdefined\hypertarget\hypertarget{sourcepointer-2}{}\fi
\noindent\textit{Source reference.} Original lemma 02S2: sources/stacks-chow/pinned-original-source.tex, source lines 3217--3217. Navigation pointer for \textbackslash{}ref\{lemma-proper-pushforward-rational-equivalence\} used inside the counted span (line 3377); the argument that consumes it is retained in the span.\par\medskip
\ifdefined\hypertarget\hypertarget{sourcepointer-3}{}\fi
\noindent\textit{Source reference.} Original lemma 02RQ: sources/stacks-chow/pinned-original-source.tex, source lines 2701--2701. Navigation pointer for \textbackslash{}ref\{lemma-rational-function\} used inside the counted span (line 3375); the argument that consumes it is retained in the span.\par\medskip
\ifdefined\hypertarget\hypertarget{sourcepointer-4}{}\fi
\noindent\textit{Source reference.} Original reference 02QL: sources/stacks-chow/pinned-original-source.tex, source lines 1372--1372. Navigation pointer for \textbackslash{}ref\{situation-setup\} used inside the counted span (line 3336); the argument that consumes it is retained in the span.\par\medskip
\begin{lemma}
\label{lemma-rational-equivalence-family}
Let $(S, \delta)$ be as in Situation \ref{situation-setup}.
Let $X$ be a scheme locally of finite type over $S$.
Let $W \subset X \times_S \mathbf{P}^1_S$ be an integral
closed subscheme of $\delta$-dimension $k + 1$.
Assume $W \not = W_0$, and $W \not = W_\infty$. Then
\begin{enumerate}
\item $W_0$, $W_\infty$ are effective Cartier divisors of $W$,
\item $W_0$, $W_\infty$ can be viewed as closed subschemes
of $X$ and
$$
[W_0]_k \sim_{rat} [W_\infty]_k,
$$
\item for any locally finite family of
integral closed subschemes
$W_i \subset X \times_S \mathbf{P}^1_S$
of $\delta$-dimension $k + 1$ with $W_i \not = (W_i)_0$ and
$W_i \not = (W_i)_\infty$ we have
$\sum ([(W_i)_0]_k - [(W_i)_\infty]_k) \sim_{rat} 0$
on $X$, and
\item for any $\alpha \in Z_k(X)$ with $\alpha \sim_{rat} 0$
there exists a locally finite family of
integral closed subschemes $W_i \subset X \times_S \mathbf{P}^1_S$
as above such that $\alpha = \sum ([(W_i)_0]_k - [(W_i)_\infty]_k)$.
\end{enumerate}
\end{lemma}

\begin{proof}
Part (1) follows from
Divisors, Lemma \ref{divisors-lemma-pullback-effective-Cartier-defined}
since the generic point
of $W$ is not mapped into $D_0$ or $D_\infty$ under the projection
$X \times_S \mathbf{P}^1_S \to \mathbf{P}^1_S$ by assumption.

\medskip\noindent
Since $X \times_S D_0 \to X$ is a closed immersion, we see that $W_0$
is isomorphic to a closed subscheme of $X$. Similarly for $W_\infty$.
The morphism $p : W \to X$ is proper as a composition of
the closed immersion $W \to X \times_S \mathbf{P}^1_S$ and the
proper morphism $X \times_S \mathbf{P}^1_S \to X$. By
Lemma \ref{lemma-rational-function} we have
$[W_0]_k \sim_{rat} [W_\infty]_k$ as cycles on $W$. Hence part (2) follows from
Lemma \ref{lemma-proper-pushforward-rational-equivalence} as clearly
$p_*[W_0]_k = [W_0]_k$ and similarly for $W_\infty$.

\medskip\noindent
The only content of statement (3) is, given parts (1) and (2), that
the collection $\{(W_i)_0, (W_i)_\infty\}$ is a locally finite collection
of closed subschemes of $X$. This is clear.

\medskip\noindent
Suppose that $\alpha \sim_{rat} 0$.
By definition this means there exist integral closed subschemes
$V_i \subset X$ of $\delta$-dimension $k + 1$ and rational
functions $f_i \in R(V_i)^*$ such that the family
$\{V_i\}_{i \in I}$ is locally finite in $X$ and such that
$\alpha = \sum (V_i \to X)_*\text{div}(f_i)$.
Let
$$
W_i \subset V_i \times_S \mathbf{P}^1_S \subset X \times_S \mathbf{P}^1_S
$$
be the closure of the graph of the rational map $f_i$ as in
Lemma \ref{lemma-rational-function}.
Then we have that $(V_i \to X)_*\text{div}(f_i)$
is equal to $[(W_i)_0]_k - [(W_i)_\infty]_k$ by that same lemma.
Hence the result is clear.
\end{proof}

\end{document}
